% !TeX root = closed-systems-from-comparison-completeness.tex
% ============================================================
\chapter{Triangle--Loop Identification as the Rigidity Bridge}
\label{ch:bridge}
% ============================================================
\section{Introduction}
Under the standing principle of closed-world admissibility
(\cref{sp:closed-world}), this chapter identifies the finite and loop
forms of \cref{sp:transport} as one and the same obstruction datum.
Closing the discrete obstruction arc, it proves that triangle failure
and loop-defect failure represent the same datum, preparing
\cref{ch:flow}; only later, through the observer-level and stitching
continuation, does this bridge feed the smooth realization chapters.
\Cref{ch:triangles,ch:descent_obstruction} isolated two a priori
different descriptions of
nontrivial morphism-level enrichment.
From \cref{ch:triangles}, the first finite obstruction to a common
diagonal gauge occurs at three points. In the triangle regime, every
pairwise restriction is diagonally alignable, but the full triple is
not. Equivalently, the corresponding stabilizer cosets have nonempty
pairwise intersections while the full triple intersection is empty.
From \cref{ch:descent_obstruction}, transport on a connected protocol,
after tree normalization, is classified by a based loop representation
into a stabilizer group. In that language, nontriviality is carried by
loop defect.
The purpose of the present chapter is to prove that, on the minimal
three-vertex subsystem carrying the first finite obstruction, these are
the same phenomenon. The triangle obstruction is not merely analogous
to loop defect. It is exactly the obstruction to descending transport
through triangle coherence.
The key point is categorical. In the free path groupoid, a directed
triangle is merely a loop. The relevant new structure appears only
after one specifies that directed triangle loops are to count as
coherent. This produces a quotient groupoid. The correct obstruction
is then not nontrivial transport \emph{on} that quotient, but rather the
failure of the original transport functor to factor through it.
The logical output of the chapter is the following equivalence on a
directed triangle subsystem:
\begin{align*}
	\text{triangle regime} & \quad\Longleftrightarrow\quad \text{failure of coherent vertex gauge}              \\
	                       & \quad\Longleftrightarrow\quad [\kappa_\ell] \neq [e]                               \\
	                       & \quad\Longleftrightarrow\quad \text{failure of descent through triangle coherence}
\end{align*}
Here \([\kappa_\ell]\) denotes the conjugacy class of the triangle
defect in the stabilizer group. The passage to conjugacy is essential:
the defect element itself depends on normalization choices, but its
triviality or nontriviality does not.
No appeal is made in this chapter to smooth realization, curvature,
filtrations, or quadratic carriers. Its role is narrower and more
rigid. It identifies the first finite diagonal obstruction with the
first obstruction to transport descent through triangle coherence.
\medskip
\noindent
\textbf{Dependence on labeled results.}
The chapter uses:
\begin{itemize}
	\item representative lift classification and tree normalization
	      (\cref{thm:lift-classification-main,thm:DO-tree-normalize}),
	\item the loop-based formulation of transport in the groupoid completion
	      (\cref{ch:descent_obstruction}),
	\item the strict minimality of the triangle regime
	      (\cref{thm:triangle}).
\end{itemize}
\medskip
\noindent
\textbf{Use in subsequent chapters.}
The chapter provides:
\begin{itemize}
	\item a canonical obstruction class carried by directed triangles,
	\item identification of the minimal triangle subsystem as the first
	      finite witness of failure of transport descent through triangle
	      coherence,
	\item the same static obstruction class carried next into the
	      observer-level irreversible descent of \cref{ch:flow}.
\end{itemize}
% ============================================================
\section{Single-object transport extracted from a connected lift}
% ============================================================
We first isolate the exact single-object transport regime determined by the
classification theory of \cref{ch:descent_obstruction}.
Let
\[
	\overline R:\mathcal I\to \Phys^{\disc}
\]
be a connected quotient protocol, and let
\[
	(x,g)
\]
be representative lift data in the sense of
\cref{thm:lift-classification-main}. Fix a base representative
\[
	x_*\in X
\]
lying over the unique value of \(\overline R\), and set
\[
	H_*:=\Stab(x_*).
\]
By \cref{thm:DO-tree-normalize}, every gauge class admits a
tree-normalized representative. We therefore fix, throughout the
chapter, a tree-normalized lift datum
\[
	(x,g)
\]
such that
\[
	x_i=x_*
	\qquad
	\text{for all }i\in \Obj(\mathcal I),
\]
and
\[
	g_\tau=e
	\qquad
	\text{for all }\tau\in \Mor(\mathcal G(\mathcal T))
\]
for a chosen spanning tree \(\mathcal T\subseteq |\mathcal I|\).
\begin{lemma}[Single-object transport values]
	\label{lem:bridge-single-object}
	For every morphism
	\[
		\alpha:i\to j
		\qquad
		\text{in }\mathcal G(\mathcal I),
	\]
	one has
	\[
		g_\alpha\in H_*.
	\]
\end{lemma}
\begin{proof}
	Since the lift is tree-normalized,
	\[
		x_i=x_j=x_*
		\qquad
		\text{for all }i,j.
	\]
	Because the transport law on the groupoid completion satisfies
	\[
		g_\alpha\cdot x_i=x_j,
	\]
	it follows that
	\[
		g_\alpha\cdot x_*=x_*.
	\]
	Thus \(g_\alpha\in \Stab(x_*)=H_*\).
\end{proof}
\begin{definition}[Finite transport subsystem]
	A \emph{finite transport subsystem} of the tree-normalized lift is a
	finite directed graph
	\[
		\mathcal N=(V,E)
	\]
	together with a graph morphism into the underlying directed graph of
	\(\mathcal G(\mathcal I)\). Equivalently, for each directed edge
	\[
		e:u\to v
		\qquad
		(u,v\in V),
	\]
	one specifies a morphism
	\[
		\alpha_e:u\to v
		\qquad
		\text{in }\mathcal G(\mathcal I).
	\]
	The associated transport assignment is the map
	\[
		\tau:E\to H_*,
		\qquad
		\tau(e):=g_{\alpha_e}.
	\]
\end{definition}
\begin{remark}[Single-object comparison regime]
	By \cref{lem:bridge-single-object}, every finite transport subsystem is
	governed by a one-object comparison groupoid with automorphism group
	\(H_*\). All representative data have been gauged to the single object
	\(x_*\), and the entire morphism locus is carried by the edge-labeling
	map
	\[
		\tau:E\to H_*.
	\]
	This is the precise single-object transport regime used below.
\end{remark}
% ============================================================
\section{Free path groupoid and triangle coherence}
% ============================================================
Let
\[
	\mathcal N=(V,E)
\]
be a finite transport subsystem with transport assignment
\[
	\tau:E\to H_*.
\]
\begin{definition}[Free path groupoid (recalled)]
	Let
	\[
		\Path^\pm(\mathcal N)
	\]
	denote the free path groupoid on \(\mathcal N\), exactly as in
	\cref{def:pathpm}: it is obtained by adjoining to each directed edge
	\[
		e:u\to v
	\]
	a formal inverse
	\[
		e^{-1}:v\to u,
	\]
	and imposing only the groupoid relations generated by identities and
	cancellation of adjacent inverse pairs.
\end{definition}
\begin{definition}[Extended transport functor]
	The edge-labeling map \(\tau:E\to H_*\) extends uniquely to a functor
	\[
		\Hol:\Path^\pm(\mathcal N)\to BH_*,
	\]
	where \(BH_*\) denotes the one-object groupoid with automorphism group
	\(H_*\). On formal inverses,
	\[
		\Hol(e^{-1})=\Hol(e)^{-1}=\tau(e)^{-1}.
	\]
\end{definition}
\begin{definition}[Directed triangle loop]
	\label{def:bridge-triangle-loop}
	A \emph{directed triangle loop} in \(\mathcal N\) is a composable loop
	of the form
	\[
		\ell=
		\bigl(
		p\xrightarrow{e_{pq}}q\xrightarrow{e_{qr}}r\xrightarrow{e_{rp}}p
		\bigr).
	\]
\end{definition}
\begin{definition}[Triangle defect]
	For a directed triangle loop
	\[
		\ell=
		\bigl(
		p\xrightarrow{e_{pq}}q\xrightarrow{e_{qr}}r\xrightarrow{e_{rp}}p
		\bigr),
	\]
	its \emph{triangle defect} is the element
	\[
		\kappa_\ell
		:=
		\tau(e_{rp})\,\tau(e_{qr})\,\tau(e_{pq})
		\in H_*.
	\]
\end{definition}
\begin{lemma}[Functorial realization of triangle defect]
	\label{lem:bridge-kappa-hol}
	For every directed triangle loop \(\ell\),
	\[
		\Hol(\ell)=\kappa_\ell.
	\]
\end{lemma}
\begin{proof}
	By definition of the functor \(\Hol\),
	\[
		\Hol(\ell)
		=
		\Hol(e_{rp})\,\Hol(e_{qr})\,\Hol(e_{pq})
		=
		\tau(e_{rp})\,\tau(e_{qr})\,\tau(e_{pq})
		=
		\kappa_\ell.
	\]
\end{proof}
\begin{definition}[Triangle coherence quotient]
	Let
	\[
		Q_\triangle:\Path^\pm(\mathcal N)\twoheadrightarrow
		\Path^\pm_\triangle(\mathcal N)
	\]
	denote the quotient of \(\Path^\pm(\mathcal N)\) by the smallest
	groupoid congruence for which every directed triangle loop becomes the
	identity at its base vertex. Thus every morphism of the form
	\[
		p\xrightarrow{e_{pq}}q\xrightarrow{e_{qr}}r\xrightarrow{e_{rp}}p
	\]
	is identified with
	\[
		\id_p,
	\]
	and the quotient is closed under composition and conjugation in the
	groupoid sense.
\end{definition}
\begin{definition}[Based loop groups]
	For each base vertex \(p\in V\), define
	\[
		\Omega_p(\mathcal N)
		:=
		\Hom_{\Path^\pm(\mathcal N)}(p,p),
		\qquad
		\Omega_p^\triangle(\mathcal N)
		:=
		\Hom_{\Path^\pm_\triangle(\mathcal N)}(p,p).
	\]
	These are groups under composition.
\end{definition}
\begin{definition}[Triangle normal subgroup]
	Fix \(p\in V\). Let
	\[
		N_{\triangle,p}\trianglelefteq \Omega_p(\mathcal N)
	\]
	be the normal subgroup generated by all conjugates
	\[
		\gamma^{-1}\ell\gamma
	\]
	where \(\ell\) is a directed triangle loop based at some vertex
	\(q\in V\), and
	\[
		\gamma:p\to q
	\]
	is any morphism in \(\Path^\pm(\mathcal N)\).
\end{definition}
\begin{theorem}[Loop-group quotient by triangle coherence]
	\label{thm:bridge-loop-quotient}
	For each base vertex \(p\in V\), the quotient functor
	\[
		Q_\triangle:\Path^\pm(\mathcal N)\twoheadrightarrow
		\Path^\pm_\triangle(\mathcal N)
	\]
	induces a surjective group homomorphism
	\[
		(Q_\triangle)_p:\Omega_p(\mathcal N)\twoheadrightarrow
		\Omega_p^\triangle(\mathcal N),
	\]
	whose kernel is exactly \(N_{\triangle,p}\). Consequently there is a
	canonical isomorphism
	\[
		\Omega_p^\triangle(\mathcal N)\cong
		\Omega_p(\mathcal N)/N_{\triangle,p}.
	\]
\end{theorem}
\begin{proof}
	Any functor between groupoids induces a homomorphism on based loop
	groups, so \(Q_\triangle\) induces
	\[
		(Q_\triangle)_p:\Omega_p(\mathcal N)\to\Omega_p^\triangle(\mathcal N).
	\]
	Surjectivity follows from surjectivity of \(Q_\triangle\) on morphisms.
	By construction of the quotient groupoid, the only new relations
	imposed beyond those already present in the free path groupoid are
	exactly the triangle coherence relations, together with all their
	consequences under composition and conjugation. On the based loop
	group at \(p\), these consequences are precisely the normal closure of
	all conjugates of directed triangle loops transported to the basepoint
	\(p\). Hence
	\[
		\ker((Q_\triangle)_p)=N_{\triangle,p}.
	\]
	The first isomorphism theorem therefore yields
	\[
		\Omega_p^\triangle(\mathcal N)\cong
		\Omega_p(\mathcal N)/N_{\triangle,p}.
	\]
\end{proof}
\begin{remark}[Why the quotient is necessary]
	Directed triangle loops are not relations in the free path groupoid.
	Before quotienting, they are merely loops. The passage to
	\(\Path^\pm_\triangle(\mathcal N)\) is exactly the imposition of
	triangle coherence. The bridge proved below compares finite triangle
	obstruction with the \emph{descent} of transport through this coherence
	quotient, not with an a priori functor already defined on the quotient.
\end{remark}
% ============================================================
\section{Gauge invariance of triangle defect}
% ============================================================
\begin{theorem}[Gauge invariance of triangle defect]
	\label{thm:triangle-invariance}
	Let \(\ell\) be a directed triangle subsystem of a finite transport
	subsystem extracted from a tree-normalized connected lift.
	Under any change of:
	\begin{enumerate}
		\item base representative \(x_*\),
		\item tree normalization,
		\item embedding of the subsystem into a larger finite transport system,
	\end{enumerate}
	the triangle defect transforms by conjugation in the stabilizer group:
	\[
		\kappa_\ell\mapsto h\,\kappa_\ell\,h^{-1}
		\qquad
		(h\in H_*).
	\]
	In particular,
	\[
		\kappa_\ell=e
		\quad\text{or}\quad
		\kappa_\ell\neq e
	\]
	is invariant.
\end{theorem}
\begin{proof}
	Under any admissible change of normalization, edge labels transform by
	vertex relabeling:
	\[
		\tau(e:u\to v)\mapsto h_v\,\tau(e)\,h_u^{-1},
	\]
	for some choice of elements \(h_v\in H_*\) at the vertices.
	Let
	\[
		\ell=
		\bigl(
		p\xrightarrow{e_{pq}}q\xrightarrow{e_{qr}}r\xrightarrow{e_{rp}}p
		\bigr).
	\]
	Then
	\[
		\kappa_\ell
		=
		\tau(e_{rp})\,\tau(e_{qr})\,\tau(e_{pq}).
	\]
	After relabeling,
	\begin{align*}
		\kappa_\ell
		 & \mapsto
		\bigl(h_p\tau(e_{rp})h_r^{-1}\bigr)
		\bigl(h_r\tau(e_{qr})h_q^{-1}\bigr)
		\bigl(h_q\tau(e_{pq})h_p^{-1}\bigr)                     \\
		 & =
		h_p\,\tau(e_{rp})\,\tau(e_{qr})\,\tau(e_{pq})\,h_p^{-1} \\
		 & =
		h_p\,\kappa_\ell\,h_p^{-1}.
	\end{align*}
	Thus the defect transforms by conjugation. Conjugation preserves the
	identity element, so triviality or nontriviality is invariant.
\end{proof}
\begin{remark}[Canonical obstruction class]
	By \cref{thm:triangle-invariance}, a directed triangle carries a
	canonical conjugacy class
	\[
		[\kappa_\ell]\in H_*/\mathrm{conj}.
	\]
	The intrinsic content of the obstruction is precisely the statement
	\[
		[\kappa_\ell]=[e]
		\quad\text{or}\quad
		[\kappa_\ell]\neq [e].
	\]
\end{remark}
% ============================================================
\section{Descent through triangle coherence}
% ============================================================
We now identify the exact obstruction to transport descent through
triangle coherence.
\begin{definition}[Descent through triangle coherence]
	The transport functor
	\[
		\Hol:\Path^\pm(\mathcal N)\to BH_*
	\]
	is said to \emph{descend through triangle coherence} if there exists a
	functor
	\[
		\Hol_\triangle:\Path^\pm_\triangle(\mathcal N)\to BH_*
	\]
	such that
	\[
		\Hol=\Hol_\triangle\circ Q_\triangle.
	\]
\end{definition}
\begin{theorem}[Descent criterion]
	\label{thm:bridge-descent}
	The transport functor \(\Hol\) descends through triangle coherence if
	and only if
	\[
		\Hol(\ell)=e
		\qquad
		\text{for every directed triangle loop }\ell.
	\]
	Equivalently, for a base vertex \(p\in V\),
	\[
		N_{\triangle,p}\subseteq \ker(\Hol_p),
	\]
	where
	\[
		\Hol_p:\Omega_p(\mathcal N)\to H_*
	\]
	is the induced based-loop homomorphism.
	In particular, for a directed triangle subsystem
	\[
		\ell=
		\bigl(
		p\xrightarrow{e_{pq}}q\xrightarrow{e_{qr}}r\xrightarrow{e_{rp}}p
		\bigr),
	\]
	one has
	\[
		\Hol(\ell)=e
		\quad\Longleftrightarrow\quad
		\kappa_\ell=e.
	\]
\end{theorem}
\begin{proof}
	Suppose first that \(\Hol\) descends through \(Q_\triangle\), say
	\[
		\Hol=\Hol_\triangle\circ Q_\triangle.
	\]
	If \(\ell\) is any directed triangle loop based at \(p\), then
	\[
		Q_\triangle(\ell)=\id_p.
	\]
	Therefore
	\[
		\Hol(\ell)=\Hol_\triangle(Q_\triangle(\ell))
		=
		\Hol_\triangle(\id_p)
		=
		e.
	\]
	Conversely, assume that
	\[
		\Hol(\ell)=e
		\qquad
		\text{for every directed triangle loop }\ell.
	\]
	Then every generator of the congruence defining
	\(\Path^\pm_\triangle(\mathcal N)\) is sent by \(\Hol\) to an identity
	morphism in \(BH_*\). Since \(BH_*\) is a groupoid, the same holds for all
	consequences of these relations under composition and conjugation.
	Hence \(\Hol\) is constant on every congruence class defining the
	triangle quotient. By the universal property of the quotient groupoid,
	there exists a unique functor
	\[
		\Hol_\triangle:\Path^\pm_\triangle(\mathcal N)\to BH_*
	\]
	such that
	\[
		\Hol=\Hol_\triangle\circ Q_\triangle.
	\]
	This proves the first equivalence.
	For the based-loop formulation, the descended functor exists exactly
	when every element of the normal subgroup generated by triangle loops is
	killed by \(\Hol_p\). By \cref{thm:bridge-loop-quotient}, this is
	equivalent to
	\[
		N_{\triangle,p}\subseteq \ker(\Hol_p).
	\]
	Finally, for a directed triangle subsystem, \cref{lem:bridge-kappa-hol}
	gives
	\[
		\Hol(\ell)=\kappa_\ell.
	\]
	Hence
	\[
		\Hol(\ell)=e
		\quad\Longleftrightarrow\quad
		\kappa_\ell=e.
	\]
\end{proof}
\begin{lemma}[Basepoint independence of descent]
	\label{lem:bridge-basepoint}
	The condition
	\[
		N_{\triangle,p}\subseteq \ker(\Hol_p)
	\]
	is independent of the choice of base vertex \(p\in V\).
\end{lemma}
\begin{proof}
	Let \(p,q\in V\), and let
	\[
		\gamma:p\to q
	\]
	be any morphism in \(\Path^\pm(\mathcal N)\). Conjugation by \(\gamma\)
	defines a group isomorphism
	\[
		c_\gamma:\Omega_p(\mathcal N)\to \Omega_q(\mathcal N),
		\qquad
		c_\gamma(\omega)=\gamma\,\omega\,\gamma^{-1}.
	\]
	This isomorphism sends conjugates of directed triangle loops based at
	\(p\) to conjugates of directed triangle loops based at \(q\). Hence
	\[
		c_\gamma(N_{\triangle,p})=N_{\triangle,q}.
	\]
	By functoriality of \(\Hol\),
	\[
		\Hol_q(c_\gamma(\omega))
		=
		\Hol(\gamma)\,\Hol_p(\omega)\,\Hol(\gamma)^{-1}.
	\]
	Therefore
	\[
		\Hol_p(\omega)=e
		\quad\Longleftrightarrow\quad
		\Hol_q(c_\gamma(\omega))=e.
	\]
	Thus
	\[
		N_{\triangle,p}\subseteq \ker(\Hol_p)
		\quad\Longleftrightarrow\quad
		N_{\triangle,q}\subseteq \ker(\Hol_q).
	\]
	So the condition is independent of the base vertex.
\end{proof}
\begin{definition}[Quotient loop defect after descent]
	Assume \(\Hol\) descends through triangle coherence. For each base
	vertex \(p\in V\), let
	\[
		\delta_p^\triangle:\Omega_p^\triangle(\mathcal N)\to H_*
	\]
	denote the based-loop homomorphism induced by the descended functor
	\[
		\Hol_\triangle:\Path^\pm_\triangle(\mathcal N)\to BH_*.
	\]
\end{definition}
\begin{theorem}[Route dependence equals loop defect on the descended quotient]
	\label{thm:bridge-route}
	Assume \(\Hol\) descends through triangle coherence, and let
	\[
		\Hol_\triangle:\Path^\pm_\triangle(\mathcal N)\to BH_*
	\]
	be the descended functor. Then for a base vertex \(p\in V\), transport
	on
	\[
		\Path^\pm_\triangle(\mathcal N)
	\]
	is endpoint-determined if and only if every loop defect in
	\[
		\Omega_p^\triangle(\mathcal N)
	\]
	is trivial.
	By \cref{lem:bridge-basepoint}, this condition is independent of the
	choice of \(p\).
\end{theorem}
\begin{proof}
	Suppose transport on \(\Path^\pm_\triangle(\mathcal N)\) is not
	endpoint-determined. Then there exist co-terminal morphisms
	\[
		\gamma_1,\gamma_2:p\to q
	\]
	such that
	\[
		\Hol_\triangle(\gamma_1)\neq \Hol_\triangle(\gamma_2).
	\]
	Since the target \(BH_*\) is a groupoid, \(\Hol_\triangle(\gamma_2)\) is
	invertible, and therefore
	\[
		\Hol_\triangle(\gamma_2^{-1}\circ\gamma_1)
		=
		\Hol_\triangle(\gamma_2)^{-1}\Hol_\triangle(\gamma_1)\neq e.
	\]
	Hence the loop
	\[
		\gamma_2^{-1}\circ\gamma_1\in \Omega_p^\triangle(\mathcal N)
	\]
	has nontrivial defect.
	Conversely, if some loop
	\[
		\ell\in \Omega_p^\triangle(\mathcal N)
	\]
	satisfies
	\[
		\delta_p^\triangle(\ell)=\Hol_\triangle(\ell)\neq e,
	\]
	then the co-terminal morphisms \(\ell\) and \(\id_p\) have different
	transport values. Thus transport is not endpoint-determined.
\end{proof}
% ============================================================
\section{Coherent vertex gauges on directed triangles}
% ============================================================
We now isolate the exact three-vertex obstruction.
\begin{definition}[Directed triangle subsystem]
	A \emph{directed triangle subsystem} of \(\mathcal N\) is a full
	three-vertex directed subgraph
	\[
		p\xrightarrow{e_{pq}}q\xrightarrow{e_{qr}}r\xrightarrow{e_{rp}}p.
	\]
	Its transport labels are
	\[
		a:=\tau(e_{pq}),
		\qquad
		b:=\tau(e_{qr}),
		\qquad
		c:=\tau(e_{rp}),
	\]
	and its triangle defect is
	\[
		\kappa_\ell=cba\in H_*.
	\]
\end{definition}
\begin{definition}[Coherent vertex gauge on a directed triangle]
	Let
	\[
		\ell=
		\bigl(
		p\xrightarrow{e_{pq}}q\xrightarrow{e_{qr}}r\xrightarrow{e_{rp}}p
		\bigr)
	\]
	be a directed triangle subsystem. A \emph{coherent vertex gauge} on
	\(\ell\) is a choice of elements
	\[
		u_p,u_q,u_r\in H_*
	\]
	such that
	\begin{align*}
		u_q & = \tau(e_{pq})\,u_p, \\
		u_r & = \tau(e_{qr})\,u_q, \\
		u_p & = \tau(e_{rp})\,u_r.
	\end{align*}
\end{definition}
\begin{lemma}[Pairwise coherence on a directed triangle]
	\label{lem:bridge-pairwise}
	Every two-edge restriction of a directed triangle subsystem admits a
	coherent vertex gauge.
\end{lemma}
\begin{proof}
	Fix a directed triangle subsystem
	\[
		p\xrightarrow{e_{pq}}q\xrightarrow{e_{qr}}r\xrightarrow{e_{rp}}p.
	\]
	For any chosen two-edge restriction, select an arbitrary source value in
	\(H_*\) at the initial vertex of the first chosen edge. The transport
	label on that edge then determines the target value uniquely, and the
	transport label on the second chosen edge determines the next value
	uniquely. Thus each two-edge restriction admits a coherent vertex
	gauge.
\end{proof}
\begin{proposition}[Global coherence on a directed triangle]
	\label{prop:bridge-global-coherence}
	Let
	\[
		\ell=
		\bigl(
		p\xrightarrow{e_{pq}}q\xrightarrow{e_{qr}}r\xrightarrow{e_{rp}}p
		\bigr)
	\]
	be a directed triangle subsystem. Then the following are equivalent.
	\begin{enumerate}[label=\textup{(\roman*)}]
		\item \(\ell\) admits a coherent vertex gauge.
		\item Its triangle defect is trivial:
		      \[
			      \kappa_\ell=\tau(e_{rp})\,\tau(e_{qr})\,\tau(e_{pq})=e.
		      \]
		\item Transport on the free path groupoid of this three-vertex subsystem
		      descends through its triangle coherence quotient.
	\end{enumerate}
\end{proposition}
\begin{proof}
	\emph{\textup{(i)}\(\Rightarrow\)\textup{(ii)}.}
	Assume a coherent vertex gauge \(u_p,u_q,u_r\) exists. Then
	\[
		u_q=\tau(e_{pq})u_p,
		\qquad
		u_r=\tau(e_{qr})u_q,
		\qquad
		u_p=\tau(e_{rp})u_r.
	\]
	Substituting successively gives
	\[
		u_p
		=
		\tau(e_{rp})\,\tau(e_{qr})\,\tau(e_{pq})\,u_p
		=
		\kappa_\ell u_p.
	\]
	Multiplying on the right by \(u_p^{-1}\) yields
	\[
		\kappa_\ell=e.
	\]
	\smallskip
	\noindent
	\emph{\textup{(ii)}\(\Rightarrow\)\textup{(i)}.}
	Assume \(\kappa_\ell=e\). Choose any \(u_p\in H_*\), and define
	\[
		u_q:=\tau(e_{pq})u_p,
		\qquad
		u_r:=\tau(e_{qr})u_q.
	\]
	Then
	\[
		\tau(e_{rp})u_r
		=
		\tau(e_{rp})\,\tau(e_{qr})\,\tau(e_{pq})\,u_p
		=
		\kappa_\ell u_p
		=
		u_p.
	\]
	Hence \(u_p,u_q,u_r\) form a coherent vertex gauge.
	\smallskip
	\noindent
	\emph{\textup{(ii)}\(\Longleftrightarrow\)\textup{(iii)}.}
	Apply \cref{thm:bridge-descent} to the present three-vertex subsystem.
	By \cref{lem:bridge-kappa-hol}, the unique directed triangle loop
	\(\ell\) satisfies
	\[
		\Hol(\ell)=\kappa_\ell.
	\]
	Therefore transport descends through triangle coherence if and only if
	\[
		\Hol(\ell)=e,
	\]
	which is equivalent to
	\[
		\kappa_\ell=e.
	\]
\end{proof}
\begin{corollary}[Directed triangles are the minimal coherent obstruction]
	\label{cor:bridge-minimal}
	For a directed triangle subsystem, every two-edge restriction is
	coherently gaugeable, while the full three-edge system fails to admit a
	coherent vertex gauge if and only if the triangle defect is nontrivial.
	Equivalently, the three-edge system is the minimal obstruction to
	descent through triangle coherence on that subsystem.
\end{corollary}
\begin{proof}
	The pairwise statement is \cref{lem:bridge-pairwise}. The full
	three-edge statement is the contrapositive of
	\cref{prop:bridge-global-coherence}.
\end{proof}
% ============================================================
\section{Single-object reformulation of the Chapter~\ref{ch:triangles} triangle regime}
% ============================================================
We now identify the exact relation between the diagonal obstruction of
\cref{ch:triangles} and coherent vertex gauges on a finite subsystem of
a tree-normalized lift.
\begin{proposition}[Diagonal alignment versus coherent vertex gauge]
	\label{prop:bridge-diagonal-gauge}
	Let \(\mathcal N=(V,E)\) be a finite transport subsystem extracted from
	a tree-normalized connected lift. Then the existence of a common
	diagonal gauge on the corresponding finite subsystem is equivalent to
	the existence of a coherent vertex gauge on \(\mathcal N\), namely an
	assignment
	\[
		u_v\in H_*
		\qquad
		(v\in V)
	\]
	such that for every directed edge
	\[
		e:u\to v
	\]
	one has
	\[
		u_v=\tau(e)\,u_u.
	\]
\end{proposition}
\begin{proof}
	Because the lift is tree-normalized, every object representative on the
	finite subsystem is equal to the single base representative \(x_*\), and
	every transport label lies in the stabilizer group \(H_*=\Stab(x_*)\). A
	choice of vertex labels \(u_v\in H_*\) therefore acts entirely within the
	single orbit of \(x_*\).
	By the diagonal-alignment criterion of \cref{ch:triangles}, common
	diagonal gaugeability is exactly the existence of a single compatible
	choice of representatives across the finite subsystem. In the present
	single-object regime, such compatibility along an edge
	\[
		e:u\to v
	\]
	means precisely that the target label is obtained from the source label
	by the transport element \(\tau(e)\):
	\[
		u_v=\tau(e)\,u_u.
	\]
	Thus the existence of a common diagonal gauge is equivalent to the
	existence of a coherent vertex gauge.
\end{proof}
\begin{corollary}[Triangle regime in the single-object transport model]
	\label{cor:bridge-triangle-regime}
	On a directed triangle subsystem extracted from a tree-normalized
	connected lift, the triangle regime of \cref{ch:triangles} is
	equivalent to pairwise coherent gaugeability together with failure of a
	coherent vertex gauge on the full triangle.
\end{corollary}
\begin{proof}
	By \cref{prop:bridge-diagonal-gauge}, common diagonal gaugeability on a
	finite subsystem is equivalent to coherent vertex gaugeability. Apply
	this to the two-edge restrictions and to the full three-edge subsystem.
\end{proof}
% ============================================================
\section{Triangle--loop identification as the rigidity bridge}
% ============================================================
We can now state the bridge theorem in exact form.
\begin{theorem}[Triangle--loop identification]
	\label{thm:bridge-main}
	Under the standing principle \cref{sp:closed-world}, let
	\[
		\ell=
		\bigl(
		p\xrightarrow{e_{pq}}q\xrightarrow{e_{qr}}r\xrightarrow{e_{rp}}p
		\bigr)
	\]
	be a directed triangle subsystem of a finite transport subsystem
	extracted from a tree-normalized connected lift. Then the following
	are equivalent.
	\begin{enumerate}[label=\textup{(\arabic*)}]
		\item The corresponding three-point subsystem is in the triangle regime
		      of \cref{ch:triangles}: every pairwise restriction admits a common
		      diagonal gauge, but the full triple does not.
		\item The directed triangle \(\ell\) admits coherent gauges on every
		      two-edge restriction, but no coherent vertex gauge on the full
		      triangle.
		\item The triangle defect is nontrivial:
		      \[
			      \kappa_\ell=\tau(e_{rp})\,\tau(e_{qr})\,\tau(e_{pq})\neq e.
		      \]
		\item The canonical conjugacy class is nontrivial:
		      \[
			      [\kappa_\ell]\neq [e]
			      \qquad
			      \text{in }H_*/\mathrm{conj}.
		      \]
		\item Transport on the free path groupoid of this three-vertex
		      subsystem does not descend through its triangle coherence quotient.
	\end{enumerate}
	If, in addition, \(\kappa_\ell=e\), so that descent exists, then the
	descended transport on the triangle quotient is endpoint-determined if
	and only if its quotient loop defects are trivial.
\end{theorem}
\begin{proof}
	\emph{\textup{(1)}\(\Longleftrightarrow\)\textup{(2)}.}
	This is \cref{cor:bridge-triangle-regime}.
	\smallskip
	\noindent
	\emph{\textup{(2)}\(\Longleftrightarrow\)\textup{(3)}.}
	By \cref{cor:bridge-minimal}, every two-edge restriction is coherently
	gaugeable, while the full triangle fails to admit a coherent vertex
	gauge if and only if \(\kappa_\ell\neq e\).
	\smallskip
	\noindent
	\emph{\textup{(3)}\(\Longleftrightarrow\)\textup{(4)}.}
	By \cref{thm:triangle-invariance}, \(\kappa_\ell\) transforms only by
	conjugation under admissible changes of normalization. Hence
	\[
		\kappa_\ell=e
		\quad\Longleftrightarrow\quad
		[\kappa_\ell]=[e],
	\]
	and therefore
	\[
		\kappa_\ell\neq e
		\quad\Longleftrightarrow\quad
		[\kappa_\ell]\neq [e].
	\]
	\smallskip
	\noindent
	\emph{\textup{(3)}\(\Longleftrightarrow\)\textup{(5)}.}
	By \cref{prop:bridge-global-coherence},
	\[
		\kappa_\ell=e
		\quad\Longleftrightarrow\quad
		\text{transport descends through triangle coherence on the subsystem}.
	\]
	Negating both sides yields
	\[
		\kappa_\ell\neq e
		\quad\Longleftrightarrow\quad
		\text{transport does not descend through triangle coherence}.
	\]
	The final statement is exactly \cref{thm:bridge-route}, applied in the
	case where descent exists.
\end{proof}
\begin{corollary}[Structural closure at the triangle boundary]
	On the minimal three-vertex subsystem carrying a finite diagonal
	obstruction, failure of transport descent through triangle coherence is
	exactly the first finite morphism-level witness of global transport
	incompatibility.
	Equivalently, the first finite witness of failure of transport descent
	through triangle coherence is triangular.
\end{corollary}
\begin{proof}
	By \cref{thm:bridge-main}, on a directed triangle subsystem the
	triangle regime of \cref{ch:triangles} is equivalent to failure of
	transport descent through triangle coherence. Since \cref{thm:triangle}
	identifies the triangle regime as the first strictly minimal finite
	witness of global incompatibility, the same subsystem is the first
	finite witness of failure of such descent.
\end{proof}
% ============================================================
\section{Interpretation}
% ============================================================
The bridge established in this chapter is exact at the level at which
it is claimed.
The finite diagonal obstruction of \cref{ch:triangles} is not merely
analogous to loop theory. On the minimal directed triangle subsystem,
it is exactly the obstruction to descending transport through triangle
coherence. Thus the first intrinsically multi-point failure of
quotient semantics is already groupoidal.
At the same time, the chapter makes no unsupported claim about graded
depth, commutator filtrations, or quadratic carriers. Those belong to
later chapters, and require their own algebraic carriers and proofs.
The role of the present chapter is narrower and more rigid: it proves
that the first finite diagonal obstruction and the first obstruction to
transport descent coincide on the minimal triangle subsystem on which
the obstruction is carried.
\begin{remark}[What is not yet used]
	Nothing in the present chapter uses smooth realization, curvature,
	connections, augmentation filtrations, or quadratic carriers. The role
	of the chapter is exactly to identify the correct finite combinatorial
	obstruction and to show that, after triangle coherence is specified,
	this obstruction is precisely failure of transport descent on the
	corresponding triangle subsystem. The immediate continuation is
	\cref{ch:flow}, where this same static obstruction is carried into
	observer-level irreversible descent; \cref{ch:stitching} then extends
	that continuation before the later curvature and interface chapters
	\cref{ch:christoffel,ch:interface} begin from it, but none of those
	later chapters is used here.
\end{remark}

\section{Conclusion}
\label{sec:bridge-conclusion}
The bridge theorem closes the discrete obstruction arc with theorem-level
precision: on the minimal triangle subsystem, finite triangular
incompatibility, nontrivial loop defect, and failure of transport descent
are equivalent manifestations of a single obstruction class.

Accordingly, \cref{ch:flow} carries this same class into the temporal
regime, where descent failure becomes the structural mechanism of
irreversible dynamics; \cref{ch:stitching} then continues that residue
through the stitching construction before the later smooth realization
chapters \cref{ch:christoffel,ch:interface}.